\chapter*{Abstract}
\addcontentsline{toc}{chapter}{Abstract}

Multi-robot systems are increasingly deployed in applications that require
persistent operation, distributed decision making, and correctness under
coupled motion, action, and information constraints. In such settings,
coordination cannot be reduced to collision avoidance, static task allocation,
or isolated path planning. The central challenge is to connect high-level task
specifications with low-level robot behavior in a manner that remains
computationally manageable and operationally meaningful when multiple robots
interact through shared resources, temporal dependence, incomplete knowledge,
limited communication, and changing mission conditions. This monograph studies
that challenge from a formal and algorithmic perspective, with emphasis on
hybrid coordination for multi-robot systems under complex temporal tasks.

The technical development begins with the nominal single-robot setting, where
robot motion is abstracted by a weighted transition system and tasks are
specified in linear temporal logic, leading to an automata-based synthesis and
hybrid execution framework. On this basis, the monograph develops two
complementary views of multi-robot coordination. The first is a bottom-up
perspective, in which robots are assigned local temporal tasks and coordination
emerges through execution-time dependence, collaborative actions, online
requests, and partial knowledge of the workspace. The second is a top-down
perspective, in which a global team task is decomposed, assigned, and refined
so that its temporal structure is preserved across heterogeneous robots and
hierarchical team organizations. The monograph then extends these foundations
to settings in which communication must be planned jointly with task
execution, and further to environments involving probabilistic task evolution,
safe-return requirements, dynamic targets, and signal-based temporal
specifications, using reactive planning, hierarchical abstraction, and
continuous trajectory optimization.

Taken together, these developments provide a unified perspective on how formal
temporal-task semantics, hybrid execution models, and scalable coordination
mechanisms can be integrated for multi-robot systems operating in complex and
evolving environments. The contribution of the monograph lies not only in the
systematic treatment of nominal, communication-constrained, and uncertain
settings within a common formal framework, but also in clarifying the
methodological links among symbolic abstraction, task decomposition, reactive
replanning, and continuous motion synthesis. At the same time, the monograph
highlights several open directions that remain central to the field, including
scalability beyond monolithic synthesis, continual coordination in
nonstationary environments, open-world reasoning under partial knowledge, and
the integration of learning-based components with meaningful formal guarantees.
